\chapter{Garbage Collection}\label{ch:garbage-collection}

Manual memory management asks every programmer to free exactly the right objects at exactly the right time.
Large programs fail that test through leaks, double frees, and dangling references.
A \term{garbage collector} reclaims unreachable objects automatically so programmers manage reachability instead of storage.
The collector must never free a live object and should reclaim dead objects promptly enough to bound memory.
This chapter builds that contract from three classic designs and one precise marking algorithm.
We trace each collection by hand so the invariants are checkable without running code.

\section{Reachability Is the Contract}\label{sec:ch33-reachability}

An object is live when the program can still reach it.
\term{Roots} are the starting set: globals, stack slots, registers, and map-described references.
An object is reachable when some path of references leads from a root to it.
Everything else is garbage, no matter how recently it was allocated.
Figure~\ref{fig:ch33-heap-graph} shows a heap with two roots, five reachable objects, and two unreachable ones.
Reachability, not age or size, decides the outcome.
Table~\ref{tab:ch33-collectors} compares the three collectors of this chapter by pause shape, throughput cost, and fragmentation behavior.
The entries are qualitative because real behavior depends on allocation rate and heap shape.
\term{Mark-sweep} marks reachable objects and sweeps the rest into free lists.
\term{Semi-space copying} moves survivors into a fresh half and frees the old half at once.
\term{Reference counting} tracks incoming edges and frees objects whose count drops to zero.
Each design upholds the same safety rule.
No live object is ever reclaimed.
They differ in when they pay, how long they pause, and whether free memory stays contiguous.

\begin{figure}[htbp]
\centering
\begin{tikzpicture}[node distance=1.0cm and 1.5cm]
\node[basic block] (r1) {root R1};
\node[basic block, below=of r1] (a) {A};
\node[basic block, below left=of a] (b) {B};
\node[basic block, below right=of a] (c) {C};
\node[basic block, right=of r1] (r2) {root R2};
\node[basic block, below=of r2] (d) {D};
\node[basic block, below=of d] (e) {E dead};
\node[basic block, left=of e] (f) {F dead};
\draw[cfg edge] (r1) -- (a);
\draw[cfg edge] (a) -- (b);
\draw[cfg edge] (a) -- (c);
\draw[cfg edge] (r2) -- (d);
\draw[cfg edge] (d) -- (a);
\end{tikzpicture}
\caption{Two roots reach A, B, C, and D through shared and direct edges, while E and F have no path from any root and are garbage.}
\label{fig:ch33-heap-graph}
\end{figure}

\begin{table}[htbp]
\centering
\caption{Collector designs uphold the same safety rule but pay at different times and fragment differently.}
\label{tab:ch33-collectors}
\begin{tabular}{llll}
\toprule
Design & Pause shape & Steady cost & Fragmentation \\
\midrule
Mark-sweep & Mark pause & Mark per live object & Free lists fragment \\
Semi-space & Copy pause & Copy per survivor & Survivors stay compact \\
Ref. counting & Small steps & Count per update & Needs cycle handling \\
\bottomrule
\end{tabular}
\end{table}

\section{Tri-Color Marking}\label{sec:ch33-tricolor}

Marking must visit every reachable object exactly once without losing track of partially scanned ones.
The \term{tri-color invariant} assigns each object a color.
White means undiscovered.
Gray means discovered but not fully scanned.
Black means scanned with all children discovered.
Collection starts with roots gray and everything else white.
Each step takes a gray object, grays its white children, and paints it black.
When no gray objects remain, black is exactly the reachable set.
Algorithm~\ref{alg:ch33-mark} states the worklist version used by mark-sweep and by concurrent markers.
It assumes the object graph is fixed during marking and that every reference is enumerable.
It requires that roots are enumerated before the loop starts.
Its output is the set of reachable objects; sweep or copy then follows.
The invariant does the heavy lifting.
No black object ever points to a white object.
That property holds initially because nothing is black.
Each step preserves it because children gray before their parent blacks.
At the end, with no gray left, no reachable object can still be white.
Mutators complicate the picture when collection runs alongside the program.
A write that installs a white pointer into a black object would hide garbage as live or, worse, hide a live object as garbage.
The standard remedy is a \term{write barrier}: a small check on pointer stores that preserves the invariant.
An insertion barrier can shade the newly stored referent; a deletion barrier can record the old referent before overwriting it; another design logs the modified object for rescanning.
These choices maintain different invariants and cannot be interchanged without changing the collector's proof.
The barrier code is short, but removing any of its cases breaks safety.

\begin{algorithm}[htbp]
\caption{Tri-color marking with an explicit worklist}\label{alg:ch33-mark}
\begin{pseudocode}
\Require valid non-null object roots $R$, stable heap $H$ with reference lists; mutator stopped
\Ensure every reachable object is marked black
\State color all objects white
\State $W \gets \emptyset$
\ForAll{$r \in R$}
\If{$r$ is white}
\State color $r$ gray and add $r$ to $W$
\EndIf
\EndFor
\While{$W \neq \emptyset$}
\State remove some $g$ from $W$
\ForAll{valid non-null object references $c$ in $g$}
\If{$c$ is white}
\State color $c$ gray and add $c$ to $W$
\EndIf
\EndFor
\State color $g$ black
\EndWhile
\State \Return all black objects
\end{pseudocode}
\end{algorithm}

\section{Worked Example: Marking and Sweeping by Hand}\label{sec:ch33-worked}

We run Algorithm~\ref{alg:ch33-mark} on the heap of Figure~\ref{fig:ch33-heap-graph}.
Roots are $R1$ and $R2$.
References are $A \to B$, $A \to C$, $D \to A$, and no edges from $B$, $C$, $E$, or $F$.
All objects start white.
Initialization grays $R1$-reachable entry $A$ and $R2$-reachable entry $D$, so $W = \{A, D\}$.
Pop $A$: children $B$ and $C$ are white, so both turn gray and $W = \{D, B, C\};$ then $A$ turns black.
Figure~\ref{fig:visual-gc-transition} freezes the state on either side of that step.
Object $D$ stays gray because it has not yet been scanned; removing it merely because another root was processed would lose pending work.

\begin{figure}[htbp]
  \centering
  \begin{tikzpicture}[node distance=11mm and 15mm]
    \node[basic block,align=center] (before)
      {Before scanning A\\white: B, C, E, F\\gray: A, D\\black: empty};
    \node[basic block,align=center,right=25mm of before] (after)
      {After scanning A\\white: E, F\\gray: D, B, C\\black: A};
    \draw[cfg edge] (before) -- node[above,font=\small] {scan A} (after);
    \node[font=\sffamily\small,align=center,below=of after] (note)
      {D remains pending.\\Scanning A discovers B and C.};
    \draw[cfg edge,dashed] (note) -- (after);
  \end{tikzpicture}
  \caption{Marking moves the scanned object to black and its newly discovered children to gray without discarding other pending roots.}
  \label{fig:visual-gc-transition}
\end{figure}


Pop $D$: child $A$ is black, so nothing changes; $D$ turns black and $W = \{B, C\}$.
Pop $B$: no children, so $B$ turns black and $W = \{C\}$.
Pop $C$: no children, so $C$ turns black and $W = \emptyset$.
The black set is $\{A, B, C, D\}$, which matches the reachable set exactly.
$E$ and $F$ stay white.
Sweep then reclaims $E$ and $F$ and leaves four survivors.
Counting confirms the work: four objects grayed, four blackened, three edges traversed, two objects freed.
A semi-space variant would copy the same four survivors and free both dead objects implicitly.
A reference-counting variant would have freed $E$ and $F$ when their last edges vanished, but it would keep any dead cycle alive.
The listing below pins the trace for tests.
\begin{irlisting}
; tri-color trace for Figure heap
; init: gray {A, D}, white {B, C, E, F}
; pop A: gray {D, B, C}, black {A}
; pop D: black {A, D}
; pop B: black {A, D, B}
; pop C: black {A, B, C, D}
; sweep: free {E, F}
\end{irlisting}
Any implementation that reports a different black set or frees a black object fails this trace.

\section{Counting, Cycles, and Barriers}\label{sec:ch33-counting}

Reference counting frees an object when its count reaches zero.
The rule is local and pauses stay short.
Its weakness is cycles: a dead ring keeps every member above zero forever.
Production counters therefore pair counting with a periodic cycle detector or a backup tracer.
Copying collectors take the opposite trade.
They compact survivors but must copy every live object on each collection.
Mark-sweep avoids copying but fragments the free space.
Fragmentation then forces size classes, coalescing, or occasional compaction.
\begin{ccwarning}[Barrier first]
A concurrent or incremental collector without a correct write barrier is a wrong collector.
Prove the barrier before tuning pause times.
\end{ccwarning}
Roots need the same rigor as heap edges.
A stack map that omits one reference slot turns a live object into garbage.
A conservative root scan that invents references turns garbage into leaks.
Test a collector implementation with adversarial heaps: shared diamonds, self-loops, two-cycles, and deep chains.
Each test asserts the exact survivor set, not merely that nothing crashed.
The \enquote{exact set} rule catches collectors that leak quietly.

\section{Generations: Most Objects Die Young}\label{sec:ch33-generations}

Measurements across many systems agree on one pattern: most allocated objects become unreachable soon after creation.
A \term{generational collector} exploits that pattern with two or more regions called generations.
New objects enter the young generation, also called the nursery.
Survivors of several nursery collections promote to the old generation.
The nursery collects often and scans only itself plus roots plus remembered old-to-young references.
A \term{remembered set} records those cross-generation pointers so the nursery need not scan the whole old heap.
The old generation collects rarely with a full traced collection.
Promotion policy controls the trade.
Promote too fast and the old generation fills with short-lived data.
Promote too slow and survivors are copied repeatedly in the nursery.
A typical rule promotes after one or two survivals, with large objects allocated directly in the old generation.
Consider a loop that builds one temporary list per iteration and keeps every tenth list.
Each iteration allocates one list and drops nine of ten.
Nursery collection reclaims the nine cheaply because it copies only the survivor.
The remembered set stays small because few old objects point into the nursery.
A generational-collector test should allocate one thousand temporaries, retain one hundred, and assert the survivor set after a nursery collection.
A second test drops all roots and asserts the nursery empties completely.
The \enquote{empty after drop} check catches remembered-set leaks that pin dead objects.

\section{Implementation Notes and Complexity}\label{sec:ch33-implementation}

Marking visits each reachable object once and traverses each of its edges once.
Tracing costs $O(|R|+L+E_L)$ where $L$ is the live-object count and $E_L$ is the number of references scanned in live objects.
The explicit initialization in the algorithm also visits all heap objects, adding $O(|H|)$ time.
Sweep scans the whole heap in $O(H)$ where $H$ is the heap size.
Copying costs $O(S)$ where $S$ is the survivor size plus forwarding work.
Updating a reference count costs $O(1)$, but a decrement to zero can trigger a cascade of reclamation; that work must be counted separately.
Memory overhead includes one color or mark bit per object, free-list links or a reserve half, and barrier metadata.
Production collectors generationalize: young objects die fast, so a small nursery collects often and survivors promote.
They run concurrently, snapshot roots atomically, and resume mutators during marking.
Each refinement preserves the tri-color invariant or states its replacement explicitly.
A common pitfall sizes the heap from a single program and declares victory.
Heap behavior changes with input size, thread count, and object shape.
Another pitfall measures only throughput and ignores pause distribution.
Interactive programs feel the longest pause, not the average.
We report both, with conditions, in the measurement chapter.
% TODO(cite): need a primary reference for write-barrier correctness arguments.

Foreign code adds roots that the ordinary managed object graph may not expose.
A retained callback or buffer export must stay reachable until the native side acknowledges that its use has ended.
Chapter~\ref{ch:ffi-ownership} develops explicit close, finalization, allocator boundaries, and cross-runtime cycles; Chapter~\ref{ch:ffi-callbacks} adds the admission-and-quiescence protocol required before reclaiming an asynchronous callback.
Lazy runtimes introduce another important root family: thunk environments and update frames.
Chapter~\ref{ch:lazy-functional-runtime} follows their references through forcing, successful update, and removal of obsolete captures without claiming that sharing always reduces residency.

\section{Exercises}\label{sec:ch33-exercises}

\begin{enumerate}
\item Run Algorithm~\ref{alg:ch33-mark} on roots $\{X\}$ with edges $X \to Y$, $Y \to X$, and isolated $Z$, list the color of each object after every pop, and give the sweep set.\label{ex:ch33-cycle-trace}
\item Give a two-object heap where reference counting without cycle detection leaks, show the counts after dropping the external edge, and state what a backup tracer would free.\label{ex:ch33-count-leak}
\item Design a write-barrier test with a black object, a white child, and one pointer store, and explain which barrier action keeps the tri-color invariant.\label{ex:ch33-barrier-test}
\end{enumerate}

\section{Further Reading}\label{sec:ch33-reading}

Appel treats storage management and collection as part of the runtime system~\cite{appel-1998-modern}.
Morgan discusses how allocation and lifetime shape optimization choices~\cite{morgan-1998-optimizing}.
Safonov addresses trustworthy reclamation when components share memory~\cite{safonov-2010-trustworthy}.
The C standard draft defines object lifetime and pointer-validity rules for C-managed memory~\cite{wg14-2011-c11}.
